\section{Proposed framework}
\label{sec:framework}

The study is organized as a prospectively specified, source-locked simulation
experiment: the bridge-state catalogues, active-mechanism scope, model-selection
policy, statistical analyses, and claim boundaries were fixed in versioned,
hash-identified source before production generation. In this paper,
``registered'' means fixed in that repository source; it does not imply deposit
in an external registry. This section gives an overview, while
Sections~\ref{sec:simulation} and~\ref{sec:data_processing} provide the
implementation details.

Three principles govern the experiment.

\paragraph{Explicit scientific blocks and common random numbers.}
The campaign contains four blocks: a two-span scour benchmark (F40-S), its
multi-damage counterpart (F40-M), a four-span multi-pier scour block (L99-S),
and its multi-damage counterpart (L99-M). Named random substreams keyed to
semantic state identifiers make all intended comparisons invariant to row
order and parallel scheduling. F40-S and F40-M share an explicit 30-state
controlled subset; their remaining states are not claimed to be paired.
L99-S and L99-M share the complete 475-state latent design. All four blocks use
one fixed rail-profile realization and vary only the registered operational
conditions between passages. Thus any paired contrast uses only states whose
semantic identities and operational draws are genuinely shared.

\paragraph{Block-local calibration.}
The architecture factorial is calibrated on F40-S using the prospectively
chosen front-bogie vertical-acceleration anchor. Frozen hyperparameters then
support the complete F40-S single- and two-channel screen. After the physical
channel pair is selected, the retained pipelines are re-optimized on that pair
in F40-S and independently in F40-M, L99-S, and L99-M. No block inherits
another block's hyperparameters for its primary result. The earlier
transport/rescue trigger is withdrawn; a separate frozen-F40-S transfer
analysis is retained only as a secondary descriptive analysis.

\paragraph{A selection/report firewall.}
States, never passages, are assigned to train, inner-validation, and sealed
outer-test partitions. Hyperparameter search, the Option-C grouped-development
out-of-fold adjudication, and the response-channel screen use development data
only. The outer test remains unopened until the architecture vectors, selected
pair, block-local hyperparameters, checkpoint policy, and comparator set are
frozen. Post-freeze variability is then estimated with a disjoint,
prospectively listed set of initialization seeds and is report-only.

\subsection{Scientific blocks and estimands}
\label{sec:estimands}

Table~\ref{tab:stages} and Figure~\ref{fig:stage_graph} summarize the campaign.
F40-S is the gating method-development block. Its dense 0--60\% grid supports
continuous regression, while the predeclared \{0, 5, 10, 20\%\} subset gives a
limited classification comparability check with Fernandes et
al.~\cite{fernandes2025early}. F40-M activates nominal bearing fixity and a
local flexural-rigidity-loss surrogate in addition to scour. L99-S estimates
three internal supports on a longer four-span bridge, and L99-M activates the
same multi-damage task there.

The design supports four distinct estimand classes:

\begin{enumerate}
  \item \textbf{F40-S method and channel selection.} Grouped-development
  evidence selects one RAW and one PAA factorial cell, and a prospectively
  specified screen over eight singles and 28 pairs selects one physical
  two-channel input. Singles are descriptive and non-selectable.
  \item \textbf{Block-local primary performance.} The retained RAW/PAA winners
  and their RAW/PAA CNN--GAP baselines receive independent selected-pair HPO in
  every block. Performance on the sealed outer test therefore estimates each
  pipeline under block-local calibration, not hyperparameter transport.
  \item \textbf{Paired damage/task comparisons.} F40-S versus F40-M is
  interpreted only on their 30 shared controlled states; L99-S versus L99-M
  uses the complete shared 475-state inventory. Each comparison jointly
  changes the expressed bearing/crack physics and, where applicable, the
  additional bearing regression heads. It is not an isolated causal effect of
  either mechanism.
  \item \textbf{Scale and transfer observations.} Comparisons between F40 and
  L99 are descriptive geometry/scale observations. The secondary frozen-F40-S
  transfer analysis measures performance under deliberately transported
  settings and cannot replace any block's independently tuned primary result.
\end{enumerate}

\begin{figure}[htbp]
\centering
\resizebox{0.88\textwidth}{!}{%
\begin{tikzpicture}[
  node distance=12mm and 26mm,
  block/.style={draw, rounded corners=2pt, minimum height=11mm,
    minimum width=30mm, align=center, font=\small\ttfamily},
  arr/.style={-{Stealth[length=2.4mm]}, thick},
  lbl/.style={font=\footnotesize, align=center}
]
\node[block] (f40s) {F40-S\\305 states};
\node[block, right=of f40s] (f40m) {F40-M\\425 states};
\draw[arr] (f40s) -- node[above,lbl] {30-state matched\\controlled subset} (f40m);
\node[block, below=of f40s] (l99s) {L99-S\\475 states};
\node[block, right=of l99s] (l99m) {L99-M\\475 states};
\draw[arr] (l99s) -- node[above,lbl] {complete 475-state\\pairing} (l99m);
\node[lbl, left=5mm of f40s] {two spans\\one scour target};
\node[lbl, left=5mm of l99s] {four spans\\three scour targets};
\node[lbl, right=5mm of f40m] {+ bearing fixity\\+ local $EI$ loss};
\node[lbl, right=5mm of l99m] {+ bearing fixity\\+ local $EI$ loss};
\end{tikzpicture}%
}
\caption{Four-block Paper-1 design. Arrows denote only authenticated semantic
state pairing; no causal or transport arrow connects the two geometries.}
\label{fig:stage_graph}
\end{figure}

\begin{table}[htbp]
\centering
\caption{Registered Paper-1 scientific blocks. Every state has 50 passages.
The profile is one fixed FRA-class-4 realization in all four blocks; speed,
temperature, and vehicle properties vary operationally.}
\label{tab:stages}
\small
\setlength{\tabcolsep}{4pt}
\begin{tabular}{llrlp{5.8cm}}
\toprule
Block & Geometry & States & Targets & Expressed bridge mechanisms \\
\midrule
\texttt{F40-S} & $2\times20$ m & 305 & support 2 & scour only; dense 61-severity $\times$ 5-replica grid \\
\texttt{F40-M} & $2\times20$ m & 425 & support 2 & scour + nominal bearing fixity + local $EI$ loss \\
\texttt{L99-S} & $4\times24.9$ m & 475 & supports 2--4 & multi-pier scour only \\
\texttt{L99-M} & $4\times24.9$ m & 475 & supports 2--4 & scour + nominal bearing fixity + local $EI$ loss \\
\bottomrule
\end{tabular}
\end{table}

Track-layer damage, wheel polygonization, and suspension damage remain present
in the single qualified source tree but are disabled by the four stage
configurations; no stripped scientific solver branch is used. Alternative
profile phases are also disabled. Their later
evaluation belongs to a separate track/train-damage study cut from the frozen
source revision; code availability is not evidence that those mechanisms were
active in Paper~1.

\subsection{Execution blocks and provenance}
\label{sec:blocks}

Each scientific block has its own authenticated generation contract,
hyperparameter studies, result manifests, and execution receipts. Within a
block, compared training arms use the same GPU model and pinned numerical
stack; physical host identity may differ only where the compatibility contract
permits it. Data generation may be distributed across independently qualified
hosts. Every reported row is bound to the generating source revision, channel
schema, protocol hashes, split identity, and artifact digests, so that a
result without a complete authenticated lineage fails closed rather than being
silently admitted.
